\chapter{Theoretical Foundations}
\label{ch:foundations}

This chapter establishes the mathematical scaffolding on which the experimental
framework rests.  We proceed from first principles: a measure-theoretic
probability space with a filtration formalizing the no-lookahead constraint,
the Kelly Criterion and its asymptotic optimality guarantee, Bayesian conjugate
updating, and the Multi-Armed Bandit formalism including UCB and Thompson
Sampling strategies.  All results stated here are standard; we provide
self-contained statements and proof sketches to anchor the computational
experiments of Chapter~\ref{ch:experiments} to rigorous theory.

% ---------------------------------------------------------------
\section{Probability Space and Filtration}
% ---------------------------------------------------------------

\begin{definition}[Filtered Probability Space]
\label{def:filtered_prob_space}
A \emph{filtered probability space} is a tuple
$(\Omega,\, \Filt,\, \{\Filt_t\}_{t \geq 0},\, \Prob)$
where $(\Omega, \Filt, \Prob)$ is a probability space and $\{\Filt_t\}$ is a
\emph{filtration}: an increasing sequence of sub-$\sigma$-algebras
$\Filt_0 \subseteq \Filt_1 \subseteq \cdots \subseteq \Filt$.
\end{definition}

In our setting the filtration is generated by observed returns.  Let
$r_1, r_2, \ldots$ denote the sequence of per-period returns.  We set
\[
    \Filt_t \;=\; \sigma(r_1,\, r_2,\, \ldots,\, r_t),
    \qquad \Filt_0 = \{\emptyset, \Omega\}.
\]
An agent's betting fraction at time $t$ must be $\Filt_{t-1}$-measurable,
meaning it depends only on information available \emph{before} $r_t$ is
revealed.  This is the formal no-lookahead constraint.

\begin{definition}[Admissible Strategy]
\label{def:admissible}
A sequence of betting fractions $\{f_t\}_{t \geq 1}$ is \emph{admissible} if
$f_t$ is $\Filt_{t-1}$-measurable for every $t$, and $f_t \in [0,1]$
almost surely.
\end{definition}

The wealth process under an admissible strategy satisfies the recursion
\begin{equation}
    W_{t} \;=\; W_{t-1}\,\bigl(1 + f_t\,r_t\bigr), \qquad W_0 = 1,
    \label{eq:wealth_recursion}
\end{equation}
which expanded telescopically gives $W_T = \prod_{t=1}^{T}(1 + f_t r_t)$.
Taking logarithms yields the central object of analysis:
\begin{equation}
    \log W_T \;=\; \sum_{t=1}^{T} \log\!\bigl(1 + f_t\,r_t\bigr).
    \label{eq:log_wealth}
\end{equation}

\begin{remark}
All numerical computations use $\log(1 + f_t r_t)$ via \texttt{numpy.log1p}
to avoid catastrophic cancellation when $f_t r_t$ is small.  Ruin is the
absorbing event $\{1 + f_t r_t \leq 0\}$, at which point $\log W_t = -\infty$.
\end{remark}

% ---------------------------------------------------------------
\section{The Kelly Criterion and Asymptotic Optimality}
% ---------------------------------------------------------------

\subsection{Binary Kelly Formula}

Consider a single bet with payoff odds $b > 0$: a fraction $f$ of wealth
returns $bf$ with probability $p$ and loses $f$ with probability $q = 1-p$.
The expected log-growth per step is
\[
    G(f) \;=\; p\log(1+bf) \;+\; q\log(1-f).
\]

\begin{proposition}[Binary Kelly Fraction]
\label{prop:binary_kelly}
The unique maximizer of $G(f)$ over $f \in [0,1]$ is
\begin{equation}
    f^* \;=\; \frac{bp - q}{b} \;=\; p - \frac{q}{b},
    \label{eq:binary_kelly}
\end{equation}
provided $bp > q$ (positive expected value); otherwise $f^* = 0$.
\end{proposition}

\begin{proof}
$G$ is strictly concave on $[0,1)$.  Computing $G'(f) = pb/(1+bf) - q/(1-f)$
and setting $G'(f) = 0$ yields \eqref{eq:binary_kelly} after clearing
denominators.  Strict concavity ($G'' < 0$) guarantees this is the unique
global maximum.
\end{proof}

\subsection{Continuous-Return Kelly Formula}

For continuous returns $r \sim \Normal(\mu,\sigma^2)$, the Kelly fraction
maximizes $\E[\log(1+fr)]$.  A second-order Taylor expansion around $f=0$
gives the tractable closed-form approximation:
\begin{equation}
    f^*_{\mathcal{N}} \;\approx\; \frac{\mu}{\sigma^2}.
    \label{eq:gaussian_kelly}
\end{equation}
This expression is used as the oracle benchmark throughout the experiments
and is clipped to $[0,1]$ to prevent over-leveraging.

\subsection{Strong Law of Large Numbers and Asymptotic Optimality}

\begin{theorem}[Kelly--Breiman Asymptotic Optimality]
\label{thm:kelly_slln}
Let $\{r_t\}_{t=1}^\infty$ be i.i.d.\ with
$\E[\log(1+f^*r)] > -\infty$, and let
$f^* = \arg\max_{f\in[0,1]}\E[\log(1+fr)]$.
\begin{enumerate}
    \item \textbf{(Asymptotic Growth Rate)}
    \[
        \frac{\log W_T^{(f^*)}}{T}
        \;\xrightarrow{\mathrm{a.s.}}\;
        \E\!\bigl[\log(1+f^*r)\bigr]
        \quad\text{as } T\to\infty.
    \]
    \item \textbf{(Dominance)} For any admissible $\{f_t\}$ with
    $\liminf_T \frac{1}{T}\sum_t \log(1+f_t r_t)
    < \E[\log(1+f^*r)]$,
    \[
        \frac{W_T^{(f^*)}}{W_T^{(f)}}
        \;\xrightarrow{\mathrm{a.s.}}\; +\infty.
    \]
\end{enumerate}
\end{theorem}

\begin{proof}[Proof sketch]
\textit{Part (1).}  By the SLLN applied to the i.i.d.\ sequence
$Z_t = \log(1+f^*r_t)$:
$\frac{1}{T}\sum_{t=1}^T Z_t \to \E[Z_1]$ a.s.
Since $\frac{1}{T}\log W_T^{(f^*)} = \frac{1}{T}\sum_t Z_t$, part (1)
follows directly.

\textit{Part (2).}  The difference in log-wealths is
$\frac{1}{T}[\log W_T^{(f^*)} - \log W_T^{(f)}]
\to \E[\log(1+f^*r)] - \liminf_T\frac{1}{T}\sum_t\log(1+f_t r_t) > 0$
a.s., so the ratio of wealths diverges.
\end{proof}

\begin{remark}[Finite-Horizon Cost]
\label{rem:finite_horizon}
Theorem~\ref{thm:kelly_slln} is asymptotic.  Over a finite horizon $T$,
Kelly betting can suffer severe drawdowns because the distribution of $W_T$
has a heavy right tail.  This motivates the CPPI risk constraint introduced
in Chapter~\ref{ch:methodology}.
\end{remark}

% ---------------------------------------------------------------
\section{Bayesian Updating and Posterior Beliefs}
% ---------------------------------------------------------------

\begin{definition}[Beta--Bernoulli Conjugate Pair]
\label{def:beta_bernoulli}
Let outcomes $X_t \in \{0,1\}$ be i.i.d.\ Bernoulli$(p)$ with unknown $p$.
The Beta prior $p \sim \mathrm{Beta}(\alpha_0,\beta_0)$ is conjugate: after
observing $n$ wins and $m$ losses the posterior is
\begin{equation}
    p \;\mid\; X_{1:n+m} \;\sim\; \mathrm{Beta}(\alpha_0+n,\;\beta_0+m).
    \label{eq:beta_posterior}
\end{equation}
The posterior mean is $\hat p = (\alpha_0+n)/(\alpha_0+\beta_0+n+m)$.
\end{definition}

\begin{proposition}[Sticky Prior Effect]
\label{prop:sticky_prior}
Suppose an agent accumulates $T_0$ observations in a stationary regime with
true win rate $p_0$, and the regime switches to $p_1 \neq p_0$ at step
$T_0+1$.  The number of additional observations required before the posterior
mean moves to within $|p_1-p_0|/2$ of $p_1$ is approximately
\begin{equation}
    T_{\mathrm{lag}} \;\approx\;
    \frac{(\alpha_0+\beta_0+T_0)\,|p_1-p_0|}{2\,|p_1-p_0|^2/(1)}
    \;=\; \frac{\alpha_0+\beta_0+T_0}{2\,|p_1-p_0|}.
    \label{eq:sticky_prior}
\end{equation}
\end{proposition}

\begin{proof}
After $T_{\mathrm{lag}}$ additional steps in the new regime, the posterior mean is
\[
    \hat p(T_0+T_{\mathrm{lag}})
    \;\approx\;
    \frac{(\alpha_0 + T_0 p_0) + T_{\mathrm{lag}} p_1}
         {\alpha_0+\beta_0+T_0+T_{\mathrm{lag}}}.
\]
Setting this equal to $(p_0+p_1)/2$ and solving for $T_{\mathrm{lag}}$ yields
\eqref{eq:sticky_prior} to leading order.
\end{proof}

\begin{remark}
Proposition~\ref{prop:sticky_prior} makes the detection-lag problem precise.
An agent with $T_0 = 200$ bull-market observations needs $O(T_0) = O(200)$
additional steps before its posterior mean shifts meaningfully toward a new
bear-market rate.  The Volatility-Augmented HMM (Chapter~\ref{ch:methodology})
bypasses this by using rolling volatility as a fast-response signal that does
not accumulate prior mass in the same way.
\end{remark}

% ---------------------------------------------------------------
\section{Multi-Armed Bandits}
% ---------------------------------------------------------------

\begin{definition}[Multi-Armed Bandit]
\label{def:mab}
A $K$-armed bandit is a tuple $(\mathcal{A},\{P_k\}_{k=1}^K)$ where
$\mathcal{A}=\{1,\ldots,K\}$ is the action set and $P_k$ is the reward
distribution of arm $k$ with mean $\mu_k$.  At each step $t$ the agent
selects arm $A_t$ and receives $R_t \sim P_{A_t}$.  The cumulative
pseudo-regret is
$\mathcal{R}_T = T\mu^* - \sum_{t=1}^T \mu_{A_t}$, $\mu^* = \max_k\mu_k$.
\end{definition}

\subsection{Upper Confidence Bound (UCB1)}

\begin{theorem}[UCB1 Regret Bound \citep{auer2002finite}]
\label{thm:ucb_regret}
UCB1, which selects $A_t = \arg\max_k\bigl[\hat\mu_k + \sqrt{2\ln t / N_k(t)}\bigr]$,
achieves
\[
    \E[\mathcal{R}_T] \;\leq\;
    \sum_{k:\,\Delta_k > 0} \frac{8\ln T}{\Delta_k}
    + \Bigl(1+\frac{\pi^2}{3}\Bigr)\sum_k \Delta_k,
\]
where $\Delta_k = \mu^*-\mu_k$.  This $O(\log T)$ bound matches the
Lai--Robbins (1985) lower bound asymptotically \citep{lai1985asymptotically}.
\end{theorem}

\subsection{Thompson Sampling}

Thompson Sampling selects arm $k$ at step $t$ by sampling
$\tilde\theta_k \sim \mathrm{Beta}(\alpha_k,\beta_k)$ and choosing
$A_t = \arg\max_k\tilde\theta_k$.

\begin{theorem}[Thompson Sampling Regret \citep{agrawal2012analysis}]
\label{thm:ts_regret}
For Bernoulli bandits, Thompson Sampling achieves
$\E[\mathcal{R}_T] = O(\sqrt{KT\ln T})$
and satisfies the Lai--Robbins instance-dependent lower bound up to constants \citep{lai1985asymptotically}.
\end{theorem}

In our framework, Thompson Sampling is combined with the Kelly formula: the
sampled $\tilde p_k$ is fed into \eqref{eq:binary_kelly} rather than used
purely for arm selection.  High posterior variance produces a wide spread of
$\tilde p_k$ draws, yielding on average smaller Kelly fractions than a
posterior-mean plug-in---an implicit, uncertainty-driven risk aversion.

\subsection{EXP3 for Adversarial Bandits}

\begin{theorem}[EXP3 Regret Bound \citep{auer2002nonstochastic}]
\label{thm:exp3_regret}
EXP3 with parameter $\gamma\in(0,1]$ achieves against any adaptive adversary:
\[
    \E[\mathcal{R}_T] \;\leq\;
    (e-1)\gamma G^* + \frac{K\ln K}{\gamma},
    \qquad
    G^* = \max_k\sum_t r_{k,t}.
\]
Optimizing $\gamma$ yields $\E[\mathcal{R}_T] = O(\sqrt{TK\ln K})$.
\end{theorem}

\begin{remark}[Fractional EXP3 Implementation]
\label{rem:exp3_variant}
Canonical EXP3 selects a single arm per round.  Our implementation uses EXP3
probability weights directly as continuous portfolio fractions scaled by a base
leverage (\S\ref{sec:agents}).  This \emph{fractional EXP3} variant does not
satisfy Theorem~\ref{thm:exp3_regret} exactly; the bound motivates its
robustness properties rather than providing a strict guarantee.  All results
label this agent ``EXP3 (fractional)'' accordingly.
\end{remark}

% ---------------------------------------------------------------
\section{Regime-Switching Environments and the HMM}
% ---------------------------------------------------------------

\begin{definition}[Hidden Markov Model]
\label{def:hmm}
A discrete-time HMM consists of:
(i) a finite hidden state space $\mathcal{S} = \{s_1,\ldots,s_M\}$;
(ii) an initial distribution $\pi$;
(iii) a transition matrix $A$ with $A_{ij} = \Prob(S_t=s_j \mid S_{t-1}=s_i)$;
(iv) emission distributions $b_j(\cdot) = p(\cdot \mid S_t=s_j)$.
Observations $\{y_t\}$ are conditionally independent given the hidden states.
\end{definition}

The optimal filter $\Prob(S_t \mid y_{1:t})$ is computed by the forward
algorithm, implemented in log-space for numerical stability:
\begin{equation}
    \ln\alpha_t(j) \;=\; \ln b_j(y_t)
    + \operatorname{logsumexp}_i\!\bigl[\ln\alpha_{t-1}(i)+\ln A_{ij}\bigr],
    \label{eq:forward_logspace}
\end{equation}
normalized by subtracting $\operatorname{logsumexp}_j(\ln\alpha_t(j))$.

\begin{remark}[Specified vs.\ Learned Emissions]
\label{rem:hmm_specified}
In Chapter~\ref{ch:methodology} the emission parameters
$(\mu_j,\sigma_j,k_j,\theta_j)$ are \emph{specified a priori} based on domain
knowledge of bull and bear return distributions, rather than estimated via
Baum--Welch EM.  With $T\approx 250$ observations, EM can overfit the emission
model; the a priori specification yields a more robust filter.  This framework
is therefore best understood as a \emph{Bayesian filter with a specified
generative model} rather than a fully learned HMM.
\end{remark}
